% !TEX root = ../../ctfp-print.tex

\lettrine[lhang=0.17]{我}{们先前已经看到，当我们在范畴} $\cat{C}$ 中固定一个对象 $a$ 时，映射 $\cat{C}(a, -)$ 是从 $\cat{C}$ 到 $\Set$ 的（共变）函子：
$$x \to \cat{C}(a, x)$$
（余域是 $\Set$ 因为 Hom 集 $\cat{C}(a, x)$ 是一个\emph{集合}。）我们将这个映射称为同态函子——我们先前也定义过它在态射上的作用。

现在让我们在这个映射中改变 $a$。我们得到一个新的映射，它将每个 $a$ 指派给同态\emph{函子} $\cat{C}(a, -)$：
$$a \to \cat{C}(a, -)$$
这是范畴 $\cat{C}$ 的对象到函子之间的映射，这些函子是函子范畴中的\emph{对象}（参见\hyperref[natural-transformations]{自然变换}一节中关于函子范畴的内容）。我们用记号 $[\cat{C}, \Set]$ 表示从 $\cat{C}$ 到 $\Set$ 的函子范畴。你可能还记得同态函子是典型的
\hyperref[representable-functors]{可表函子}。

每当我们有在两个范畴之间的对象映射时，自然会问这种映射是否也是一个函子。换句话说，我们能否将一个范畴中的态射提升为另一个范畴中的态射。$\cat{C}$ 中的态射就是 $\cat{C}(a, b)$ 的元素，而函子范畴 $[\cat{C}, \Set]$ 中的态射是自然变换。因此我们需要寻找态射到自然变换的映射。

让我们看看能否找到对应于态射 $f \Colon a \to b$ 的自然变换。首先，观察 $a$ 和 $b$ 被映射到什么。它们被映射到两个函子：$\cat{C}(a, -)$ 和 $\cat{C}(b, -)$。我们需要在这两个函子之间建立自然变换。

这里有个技巧：我们应用约恩达引理：
$$[\cat{C}, \Set](\cat{C}(a, -), F) \cong F a$$
并将一般的 $F$ 替换为同态函子 $\cat{C}(b, -)$。我们得到：
$$[\cat{C}, \Set](\cat{C}(a, -), \cat{C}(b, -)) \cong \cat{C}(b, a)$$
\begin{figure}[H]
  \centering
  \includegraphics[width=0.6\textwidth]{images/yoneda-embedding.jpg}
\end{figure}

\noindent
这正是我们所寻找的两个同态函子之间的自然变换，但有个微妙之处：我们得到了自然变换与态射之间的对应关系——即 $\cat{C}(b, a)$ 的元素——但方向是「相反」的。但这没问题；这只是意味着我们正在研究的函子是反变的。

\begin{figure}[H]
  \centering
  \includegraphics[width=0.65\textwidth]{images/yoneda-embedding-2.jpg}
\end{figure}

\noindent
实际上，我们获得的比预期更多。从 $\cat{C}$ 到 $[\cat{C}, \Set]$ 的映射不仅是一个反变函子——它还是一个\emph{完全忠实}的函子。满性与忠实性描述了函子如何映射Hom集。

一个\emph{忠实}函子在Hom集上是\newterm{单射}的，即它将不同的态射映射为不同的态射。换句话说，它不会合并它们。

一个\emph{全}函子在Hom集上是\newterm{满射}的，即它将一个Hom集\emph{满射}到另一个Hom集，完全覆盖后者。

一个完全忠实函子 $F$ 在Hom集上是\newterm{双射}的——这两个集合的所有元素一一对应。对于源范畴 $\cat{C}$ 中的任意对象对 $a$ 和 $b$，都存在 $\cat{C}(a, b)$ 与 $\cat{D}(F a, F b)$ 之间的双射，其中 $\cat{D}$ 是 $F$ 的目标范畴（在我们的案例中是函子范畴 $[\cat{C}, \Set]$）。请注意，这并不表示 $F$ 在\emph{对象}上是双射。目标范畴 $\cat{D}$ 中可能存在不在 $F$ 像中的对象，而且我们无法对这些对象的Hom集做出任何断言。

\section{嵌入}

我们刚刚描述的（反变）函子，即将 $\cat{C}$ 中的对象映射到 $[\cat{C}, \Set]$ 中的函子：
$$a \to \cat{C}(a, -)$$
定义了\newterm{约恩达嵌入}。它将范畴 $\cat{C}$（严格来说是范畴 $\cat{C}^\mathit{op}$，因为反变性）嵌入到函子范畴 $[\cat{C}, \Set]$ 内部。它不仅将对象映射到函子，还忠实地保留了它们之间的所有联系。

这是一个非常有用的结果，因为数学家对函子范畴（尤其是以 $\Set$ 为余域的函子）了解甚多。通过将任意范畴 $\cat{C}$ 嵌入函子范畴，我们可以获得大量洞察。

当然，约恩达嵌入也有对偶版本，有时被称为协约恩达嵌入。注意到我们也可以通过固定每个Hom集 $\cat{C}(-, a)$ 的目标对象（而非源对象）开始研究。这会给出一个反变同态函子。从 $\cat{C}$ 到 $\Set$ 的反变函子就是我们熟悉的预层（例如参见
\hyperref[limits-and-colimits]{极限与余极限}）。协约恩达嵌入定义了范畴 $\cat{C}$ 到预层范畴的嵌入。它在态射上的作用由下式给出：
$$[\cat{C}^\mathit{op}, \Set](\cat{C}(-, a), \cat{C}(-, b)) \cong \cat{C}(a, b)$$
同样，数学家对预层范畴了解甚多，因此能够将任意范畴嵌入其中是一个重大优势。

\section{在Haskell中的应用}

在Haskell中，Yoneda嵌入可表示为读者函子(reader functor)之间的自然变换与反向函数之间的一个同构：

\begin{snipv}
forall x. (a -> x) -> (b -> x) \ensuremath{\cong} b -> a
\end{snipv}
（注意：读者函子等价于\code{((->) a)}）

该恒等式的左侧是一个多态函数，给定一个从\code{a}到\code{x}的函数和类型为\code{b}的值时，可以生成类型为\code{x}的值（此处进行了去柯里化处理——省略了围绕函数\code{b -> x}的括号）。对于所有\code{x}都能实现此操作的唯一方式，是我们的函数内部已知晓如何将\code{b}转换为\code{a}，它必须隐式持有某个转换函数\code{b -> a}。

给定这样的转换器\code{btoa}，我们可以定义左侧函数（记作\code{fromY}）如下：

\src{snippet01}
反之，若已有函数\code{fromY}，我们可以通过传入恒等函数来恢复转换器：

\src{snippet02}
这就在类型为\code{fromY}和\code{btoa}的函数之间建立了双射关系。

对该同构的另一种解读视角是：它是从\code{b}到\code{a}函数的一种\acronym{CPS}（延续传递风格）编码。其中参数\code{a -> x}是一个延续（处理程序），结果则是从\code{b}到\code{x}的函数。当以类型\code{b}的值调用时，该函数会执行预先组合了待编码函数的延续。

Yoneda嵌入还解释了Haskell中数据结构的一些替代表示形式。特别是，它为\code{Control.Lens}库中的透镜(lens)提供了\urlref{https://bartoszmilewski.com/2015/07/13/from-lenses-to-yoneda-embedding/}{一种非常有用的表示法}。

\section{预序示例}

此示例由Robert Harper提出。这是将约恩达嵌入应用于由预序定义的范畴的例子。预序(preorder)是其元素间具有排序关系的集合，传统上用$\leqslant$（小于等于）表示。"预序"中的"预"字源于我们仅要求该关系满足传递性和自反性，而不必满足反对称性（因此可能存在环路）。

带有预序关系的集合可生成一个范畴。其对象是该集合的元素。当且仅当$a \leqslant b$时，存在从对象$a$到$b$的唯一态射；否则不存在此类态射。任意两个对象之间至多存在一个态射。因此此类范畴中的Hom集要么是空集，要么是单元素集。这种范畴被称为\emph{薄范畴}(thin)。

容易验证此构造确实构成范畴：若$a \leqslant b$且$b \leqslant c$，则$a \leqslant c$，故态射可复合；且复合运算满足结合律。由于每个元素都满足自反性（即$a \leqslant a$），我们还拥有单位态射。

现可对预序范畴应用共变约恩达嵌入。特别地，我们关注其在态射上的作用：
$$[\cat{C}, \Set](\cat{C}(-, a), \cat{C}(-, b)) \cong \cat{C}(a, b)$$
右侧Hom集非空当且仅当$a \leqslant b$成立——此时它是一个单元素集。因此，若$a \leqslant b$存在，则左侧存在唯一自然变换；否则不存在自然变换。

那么预序范畴中同态函子间的自然变换是什么？它应是连接各集合$\cat{C}(-, a)$与$\cat{C}(-, b)$的函数族。在预序中，这些集合要么为空，要么为单元素集。让我们考察可用的函数类型。

存在空集到自身的恒等函数，存在从空集到单元素集的\code{absurd}函数（该函数实际上不执行任何操作，因为它只需定义在空集元素上，而空集无元素），也存在单元素集到自身的恒等函数。唯一禁止的情况是从单元素集到空集的映射（该函数作用于唯一元素时会得到什么值呢？）。

因此我们的自然变换永远不会连接单元素Hom集到空Hom集。换句话说，若$x \leqslant a$（单元素Hom集$\cat{C}(x, a)$），则$\cat{C}(x, b)$不能为空集。非空的$\cat{C}(x, b)$意味着$x$小于等于$b$。因此，该自然变换的存在性要求：对所有$x$，若$x \leqslant a$则必有$x \leqslant b$。
$$\text{对所有 } x,\ x \leqslant a \Rightarrow x \leqslant b$$
另一方面，共变约恩达引理告诉我们此自然变换的存在性等价于$\cat{C}(a, b)$非空，即$a \leqslant b$。综上可得：
$$a \leqslant b \text{ 当且仅当 对所有 } x,\ x \leqslant a \Rightarrow x \leqslant b$$
我们本可以直接得出此结论。直观地说，若$a \leqslant b$，则所有低于$a$的元素必然也低于$b$。反之，当在右侧用$a$替代$x$时，可推出$a \leqslant b$。但您必须承认，通过约恩达嵌入获得此结果的过程更具启发性。

\section{自然性}

Yoneda引理建立了自然变换集合与$\Set$范畴中对象之间的同构。自然变换是函子范畴$[\cat{C}, \Set]$中的态射。任意两个函子之间的自然变换集合即该范畴中的一个同态集合。Yoneda引理给出的同构是：
$$[\cat{C}, \Set](\cat{C}(a, -), F) \cong F a$$
这个同构在$F$和$a$中都是自然的。换句话说，它在$(F, a)$这对对象上是自然的，其中$(F, a)$取自积范畴$[\cat{C}, \Set] \times \cat{C}$。请注意我们现在将$F$视为函子范畴中的\newterm{对象}。

让我们稍作思考这表示什么含义。自然同构是两个函子间的可逆\emph{自然变换}。事实上，我们同构的右侧是一个函子。它是从$[\cat{C}, \Set] \times \cat{C}$到$\Set$的函子。其作用在$(F, a)$这对对象上时得到一个集合——即函子$F$在对象$a$处的取值。这被称为评估函子。

左侧也是一个函子，它将$(F, a)$映射到自然变换集合$[\cat{C}, \Set](\cat{C}(a, -), F)$。

为了证明这些确实是函子，我们还需要定义它们在态射上的作用。但一对$(F, a)$和$(G, b)$之间的态射是什么？这是由两个态射组成的对$(\Phi, f)$：第一个是函子间的态射——即自然变换；第二个是$\cat{C}$中的常规态射。

评估函子将这对$(\Phi, f)$映射到两个集合$F a$和$G b$之间的函数。我们可以很容易地通过$\Phi$在$a$处的分量（将$F a$映射到$G a$）和$G$提升的态射$f$构造这样的函数：
$$(G f) \circ \Phi_a$$
注意由于$\Phi$的自然性，这等价于：
$$\Phi_b \circ (F f)$$
我不会证明整个同构的自然性——一旦你确定了函子的具体形式，证明过程就相当机械。这源于我们的同构是由函子和自然变换构建而成的事实。这种结构天然保证了其必然成立。

\section{挑战}

\begin{enumerate}
  \tightlist
  \item
        将余约纳嵌入（co-Yoneda embedding）表达为Haskell形式。
  \item
        证明我们在\code{fromY}与\code{btoa}之间建立的双射是一个同构（这两个映射互为逆映射）。
  \item
        推导幺半群（monoid）的约纳嵌入。对应于幺半群单对象的函子是什么？对应于幺半群态射的自然变换是什么？
  \item
        协变（covariant）约纳嵌入在预序（preorders）上的应用是什么？（问题由Gershom Bazerman提出）
  \item
        约纳嵌入可以将任意函子范畴 $[\cat{C}, \cat{D}]$ 嵌入到函子范畴 $[[\cat{C}, \cat{D}], \Set]$ 中。弄清楚它在态射（在这种情况下是自然变换）上的作用方式。
\end{enumerate}
